% pracovni_list.tex -- Lekce 07: Seznamy a hudba
\newcommand{\lekce}{Lekce 07 -- Pracovní list}
% preamble-pl.tex -- sdílená preamble pro pracovní listy
% Includuje se z ucitel/lekceXX/pracovni_list.tex jako:
%   % preamble-pl.tex -- sdílená preamble pro pracovní listy
% Includuje se z ucitel/lekceXX/pracovni_list.tex jako:
%   % preamble-pl.tex -- sdílená preamble pro pracovní listy
% Includuje se z ucitel/lekceXX/pracovni_list.tex jako:
%   \input{../spolecne/preamble-pl.tex}

\documentclass[a4paper,12pt]{article}

% --- Jazyk a font ---
\usepackage{polyglossia}
\setmainlanguage{czech}
\usepackage{fontspec}

% --- Okraje ---
\usepackage[top=20mm, bottom=20mm, left=22mm, right=22mm]{geometry}

% --- Česky korektní uvozovky ---
\usepackage{csquotes}
\newcommand{\uv}[1]{\enquote{#1}}

% --- Typografie ---
\usepackage{microtype}
\usepackage{parskip}          % odstavce odděleny mezerou, ne odsazením

% --- Barvy a grafika ---
\usepackage[table]{xcolor}
\definecolor{microbit-blue}{HTML}{1E4D8C}
\definecolor{code-bg}{HTML}{F5F5F5}
\definecolor{task-bg}{HTML}{EEF4FF}

% --- Tabulky ---
\usepackage{booktabs}
\usepackage{multirow}

% --- Zdrojový kód (Python) ---
\usepackage{listings}
\lstset{
  language=Python,
  basicstyle=\ttfamily\small,
  backgroundcolor=\color{code-bg},
  frame=single,
  framerule=0pt,
  rulecolor=\color{gray!30},
  xleftmargin=6pt,
  xrightmargin=6pt,
  breaklines=true,
  showstringspaces=false,
  keywordstyle=\color{microbit-blue}\bfseries,
  stringstyle=\color{teal},
  commentstyle=\color{gray}\itshape,
  literate=
    {á}{{\'a}}1 {é}{{\'e}}1 {í}{{\'i}}1 {ó}{{\'o}}1 {ú}{{\'u}}1
    {ů}{{\r{u}}}1 {č}{{\v{c}}}1 {š}{{\v{s}}}1 {ž}{{\v{z}}}1
    {Á}{{\'A}}1 {É}{{\'E}}1 {Í}{{\'I}}1 {Ó}{{\'O}}1 {Ú}{{\'U}}1
    {Č}{{\v{C}}}1 {Š}{{\v{S}}}1 {Ž}{{\v{Z}}}1 {ň}{{\v{n}}}1 {ř}{{\v{r}}}1,
}

% --- Zadání úkolů ---
\usepackage[most]{tcolorbox}
\tcbuselibrary{breakable}

\newcounter{ukol}[section]
\newenvironment{ukol}[1][]{%
  \stepcounter{ukol}%
  \begin{tcolorbox}[
    colback=task-bg,
    colframe=microbit-blue,
    fonttitle=\bfseries,
    title={Úkol~\theukol\ifx\relax#1\relax\else{: #1}\fi},
    breakable,
    left=6pt, right=6pt, top=4pt, bottom=4pt
  ]%
}{%
  \end{tcolorbox}%
}

% Inline kód — underscore package zajistí, ze _ funguje v texttt
\usepackage{underscore}
\newcommand{\pyinline}[1]{\texttt{#1}}
% Pojem (zvýraznění termínu) — stejné jako v preamble-prez.tex
\newcommand{\pojem}[1]{\textbf{#1}}

% Místo pro odpověď studenta
\newcommand{\odpoved}[1][3cm]{%
  \vspace{#1}%
}

% --- Záhlaví / zápatí ---
\usepackage{fancyhdr}
\pagestyle{fancy}
\fancyhf{}
\renewcommand{\headrulewidth}{0.4pt}
\fancyhead[L]{\small\textcolor{microbit-blue}{\textbf{Úvod do Pythonu na Micro:bitu}}}
\fancyhead[R]{\small\textcolor{gray}{\lekce}}
\fancyfoot[C]{\small\thepage}

% --- Ostatní ---
\usepackage{enumitem}
\usepackage{graphicx}
\usepackage{hyperref}
\hypersetup{
  colorlinks=true,
  linkcolor=microbit-blue,
  urlcolor=microbit-blue,
}


\documentclass[a4paper,12pt]{article}

% --- Jazyk a font ---
\usepackage{polyglossia}
\setmainlanguage{czech}
\usepackage{fontspec}

% --- Okraje ---
\usepackage[top=20mm, bottom=20mm, left=22mm, right=22mm]{geometry}

% --- Česky korektní uvozovky ---
\usepackage{csquotes}
\newcommand{\uv}[1]{\enquote{#1}}

% --- Typografie ---
\usepackage{microtype}
\usepackage{parskip}          % odstavce odděleny mezerou, ne odsazením

% --- Barvy a grafika ---
\usepackage[table]{xcolor}
\definecolor{microbit-blue}{HTML}{1E4D8C}
\definecolor{code-bg}{HTML}{F5F5F5}
\definecolor{task-bg}{HTML}{EEF4FF}

% --- Tabulky ---
\usepackage{booktabs}
\usepackage{multirow}

% --- Zdrojový kód (Python) ---
\usepackage{listings}
\lstset{
  language=Python,
  basicstyle=\ttfamily\small,
  backgroundcolor=\color{code-bg},
  frame=single,
  framerule=0pt,
  rulecolor=\color{gray!30},
  xleftmargin=6pt,
  xrightmargin=6pt,
  breaklines=true,
  showstringspaces=false,
  keywordstyle=\color{microbit-blue}\bfseries,
  stringstyle=\color{teal},
  commentstyle=\color{gray}\itshape,
  literate=
    {á}{{\'a}}1 {é}{{\'e}}1 {í}{{\'i}}1 {ó}{{\'o}}1 {ú}{{\'u}}1
    {ů}{{\r{u}}}1 {č}{{\v{c}}}1 {š}{{\v{s}}}1 {ž}{{\v{z}}}1
    {Á}{{\'A}}1 {É}{{\'E}}1 {Í}{{\'I}}1 {Ó}{{\'O}}1 {Ú}{{\'U}}1
    {Č}{{\v{C}}}1 {Š}{{\v{S}}}1 {Ž}{{\v{Z}}}1 {ň}{{\v{n}}}1 {ř}{{\v{r}}}1,
}

% --- Zadání úkolů ---
\usepackage[most]{tcolorbox}
\tcbuselibrary{breakable}

\newcounter{ukol}[section]
\newenvironment{ukol}[1][]{%
  \stepcounter{ukol}%
  \begin{tcolorbox}[
    colback=task-bg,
    colframe=microbit-blue,
    fonttitle=\bfseries,
    title={Úkol~\theukol\ifx\relax#1\relax\else{: #1}\fi},
    breakable,
    left=6pt, right=6pt, top=4pt, bottom=4pt
  ]%
}{%
  \end{tcolorbox}%
}

% Inline kód — underscore package zajistí, ze _ funguje v texttt
\usepackage{underscore}
\newcommand{\pyinline}[1]{\texttt{#1}}
% Pojem (zvýraznění termínu) — stejné jako v preamble-prez.tex
\newcommand{\pojem}[1]{\textbf{#1}}

% Místo pro odpověď studenta
\newcommand{\odpoved}[1][3cm]{%
  \vspace{#1}%
}

% --- Záhlaví / zápatí ---
\usepackage{fancyhdr}
\pagestyle{fancy}
\fancyhf{}
\renewcommand{\headrulewidth}{0.4pt}
\fancyhead[L]{\small\textcolor{microbit-blue}{\textbf{Úvod do Pythonu na Micro:bitu}}}
\fancyhead[R]{\small\textcolor{gray}{\lekce}}
\fancyfoot[C]{\small\thepage}

% --- Ostatní ---
\usepackage{enumitem}
\usepackage{graphicx}
\usepackage{hyperref}
\hypersetup{
  colorlinks=true,
  linkcolor=microbit-blue,
  urlcolor=microbit-blue,
}


\documentclass[a4paper,12pt]{article}

% --- Jazyk a font ---
\usepackage{polyglossia}
\setmainlanguage{czech}
\usepackage{fontspec}

% --- Okraje ---
\usepackage[top=20mm, bottom=20mm, left=22mm, right=22mm]{geometry}

% --- Česky korektní uvozovky ---
\usepackage{csquotes}
\newcommand{\uv}[1]{\enquote{#1}}

% --- Typografie ---
\usepackage{microtype}
\usepackage{parskip}          % odstavce odděleny mezerou, ne odsazením

% --- Barvy a grafika ---
\usepackage[table]{xcolor}
\definecolor{microbit-blue}{HTML}{1E4D8C}
\definecolor{code-bg}{HTML}{F5F5F5}
\definecolor{task-bg}{HTML}{EEF4FF}

% --- Tabulky ---
\usepackage{booktabs}
\usepackage{multirow}

% --- Zdrojový kód (Python) ---
\usepackage{listings}
\lstset{
  language=Python,
  basicstyle=\ttfamily\small,
  backgroundcolor=\color{code-bg},
  frame=single,
  framerule=0pt,
  rulecolor=\color{gray!30},
  xleftmargin=6pt,
  xrightmargin=6pt,
  breaklines=true,
  showstringspaces=false,
  keywordstyle=\color{microbit-blue}\bfseries,
  stringstyle=\color{teal},
  commentstyle=\color{gray}\itshape,
  literate=
    {á}{{\'a}}1 {é}{{\'e}}1 {í}{{\'i}}1 {ó}{{\'o}}1 {ú}{{\'u}}1
    {ů}{{\r{u}}}1 {č}{{\v{c}}}1 {š}{{\v{s}}}1 {ž}{{\v{z}}}1
    {Á}{{\'A}}1 {É}{{\'E}}1 {Í}{{\'I}}1 {Ó}{{\'O}}1 {Ú}{{\'U}}1
    {Č}{{\v{C}}}1 {Š}{{\v{S}}}1 {Ž}{{\v{Z}}}1 {ň}{{\v{n}}}1 {ř}{{\v{r}}}1,
}

% --- Zadání úkolů ---
\usepackage[most]{tcolorbox}
\tcbuselibrary{breakable}

\newcounter{ukol}[section]
\newenvironment{ukol}[1][]{%
  \stepcounter{ukol}%
  \begin{tcolorbox}[
    colback=task-bg,
    colframe=microbit-blue,
    fonttitle=\bfseries,
    title={Úkol~\theukol\ifx\relax#1\relax\else{: #1}\fi},
    breakable,
    left=6pt, right=6pt, top=4pt, bottom=4pt
  ]%
}{%
  \end{tcolorbox}%
}

% Inline kód — underscore package zajistí, ze _ funguje v texttt
\usepackage{underscore}
\newcommand{\pyinline}[1]{\texttt{#1}}
% Pojem (zvýraznění termínu) — stejné jako v preamble-prez.tex
\newcommand{\pojem}[1]{\textbf{#1}}

% Místo pro odpověď studenta
\newcommand{\odpoved}[1][3cm]{%
  \vspace{#1}%
}

% --- Záhlaví / zápatí ---
\usepackage{fancyhdr}
\pagestyle{fancy}
\fancyhf{}
\renewcommand{\headrulewidth}{0.4pt}
\fancyhead[L]{\small\textcolor{microbit-blue}{\textbf{Úvod do Pythonu na Micro:bitu}}}
\fancyhead[R]{\small\textcolor{gray}{\lekce}}
\fancyfoot[C]{\small\thepage}

% --- Ostatní ---
\usepackage{enumitem}
\usepackage{graphicx}
\usepackage{hyperref}
\hypersetup{
  colorlinks=true,
  linkcolor=microbit-blue,
  urlcolor=microbit-blue,
}


\begin{document}

\begin{center}
  {\Large\bfseries\textcolor{microbit-blue}{Lekce 07: Seznamy a hudba}}\\[4pt]
  {\large Pracovní list}
\end{center}

\medskip\noindent
Dosud každá proměnná uchovávala \textbf{jednu} hodnotu.
Dnes přichází nový koncept: \pojem{seznam} — proměnná, která drží
\textbf{více hodnot najednou}.
K~jednotlivým prvkům se dostaneme přes \pojem{index} — číslo pozice (od~0).

\medskip\noindent
A uvidíme, že melodie je přesně tohle: seznam not, jedna za druhou.

% ------------------------------------------------------------------
\begin{ukol}[První melodie]

Micro:bit přehraje předdefinovanou melodii po stisku tlačítka.

\begin{lstlisting}
from microbit import *
import music

while True:
    if button_a.was_pressed():
        music.play(music.BIRTHDAY)
    sleep(50)
\end{lstlisting}

\begin{enumerate}
  \item Zadej program a stiskni tlačítko A — uslyšíš melodii?
  \item Zkus zaměnit \pyinline{music.BIRTHDAY} za jinou:
        \pyinline{music.NYAN}, \pyinline{music.FUNK}, \pyinline{music.DADADADUM}.
  \item Přidej tlačítko B pro druhou melodii (použij \pyinline{elif}).
  \item \pyinline{music.play()} hraje blokujícím způsobem — program čeká, dokud melodie neskončí.
        Zkus stisknout A při přehrávání.
        Jak bys melodii přerušil? \textit{(Tip: \pyinline{music.stop()})}
\end{enumerate}

\odpoved[2cm]

\begin{tcolorbox}[colback=code-bg, colframe=gray!40,
                  fonttitle=\small\bfseries, title=Seznamy a hudba]
  \small
  \begin{itemize}[leftmargin=*]
    \item \pyinline{seznam[0]} — první prvek; \pyinline{seznam[-1]} — poslední prvek.
    \item \pyinline{len(seznam)} — počet prvků v seznamu.
    \item \pyinline{import music} — nutný navíc vedle \pyinline{from microbit import *}.
    \item \pyinline{music.play(music.NYAN)} — přehraje předdefinovanou melodii.
    \item \pyinline{music.play(seznam\_not)} — přehraje vlastní melodii ze seznamu.
  \end{itemize}
\end{tcolorbox}

\end{ukol}

% ------------------------------------------------------------------
\begin{ukol}[Vlastní melodie]

Vlastní melodii zapíšeme jako seznam řetězců — každá nota obsahuje notové jméno,
číslo oktávy a délku oddělenou dvojtečkou.
Délky: \pyinline{1}~=~celá, \pyinline{2}~=~půlová, \pyinline{4}~=~čtvrťová, \pyinline{8}~=~osminová.
\textbf{Pozor:} čím větší číslo, tím \emph{kratší} nota — jako když krájíš pizzu
na více dílů, každý díl je menší.

\begin{lstlisting}
import music

melodie = [
    "C4:4", "C4:4", "G4:4", "G4:4",
    "A4:4", "A4:4", "G4:2",
    "R:2",
    "F4:4", "F4:4", "E4:4", "E4:4",
    "D4:4", "D4:4", "C4:2"
]

music.play(melodie)
\end{lstlisting}

\begin{enumerate}
  \item Zadej melodii a poslechni si, co hraje. Poznáš ji?
  \item Přidej do melodie pomlku \pyinline{"R:4"} uprostřed — kde přesně a jak ovlivní zvuk?
  \item Vymysli nebo najdi vlastní jednoduchou melodii (aspoň 8 not) a zapiš ji jako seznam.
        Ověř, že hraje správně.
\end{enumerate}

\odpoved[3cm]

\end{ukol}

% ------------------------------------------------------------------
\begin{ukol}[Jukebox]

Propoj seznamy, hudbu a tlačítka do jednoduchého jukeboxu.

\medskip
Podmínky:
\begin{itemize}
  \item program nabízí alespoň \textbf{2 různé melodie} — jedna na tlačítko A,
        druhá na tlačítko B
  \item alespoň jedna melodie musí být \textbf{vlastní seznam not}
        (ne jen předdefinovaná)
  \item před přehráním zobraz na displeji číslo nebo písmeno (A/B),
        aby bylo jasné, která melodie hraje
\end{itemize}

\odpoved[5cm]

\noindent Výsledný program ulož a připrav k odevzdání do Google Classroom.

\end{ukol}

% ------------------------------------------------------------------
\begin{bonusukol}[Generátor náhodné melodie]

Každé zatřesení vygeneruje novou, jedinečnou melodii z předem daných not.
Propojení \pyinline{random.choice()} (lekce~5) se seznamy a \pyinline{append()}:

\begin{lstlisting}
from microbit import *
import music
import random

noty    = ["C4", "D4", "E4", "F4", "G4", "A4", "B4", "C5"]
delky   = ["4", "4", "4", "2", "8"]   # cetnost ovlivni rytmus

while True:
    if accelerometer.was_gesture("shake"):
        melodie = []
        for i in range(8):
            nota  = random.choice(noty)
            delka = random.choice(delky)
            melodie.append(nota + ":" + delka)
        music.play(melodie)
    sleep(100)
\end{lstlisting}

\begin{enumerate}
  \item Zadej program a zatřes — každá melodie je jiná!
  \item Co dělá \pyinline{melodie.append(...)}? Proč začínáme s prázdným seznamem?
  \item Uprav seznam \pyinline{noty} — vynech F a B (pentatonická stupnice).
        Zní melodie příjemněji? Proč?
  \item \textbf{Výzva:} Přidej do \pyinline{delky} více čtvrťových not
        (\pyinline{"4"}) — jak to ovlivní charakter melodie?
\end{enumerate}

\odpoved[2.5cm]

\end{bonusukol}

\end{document}
